% !TEX root = ../main.tex
\corrigesChapitre{2}{Absorbance et dosage par étalonnage}
\corrigesSouples

\begin{blocLivre}{cor-2-1}
\begin{corrige}{1}{Spectre et couleur}
\demarche{La couleur perçue est la complémentaire de la couleur la plus absorbée, lue sur le cercle chromatique.}
\begin{questions}
\item Solution 1 : $\res{\lambda_{\max}\approx\qty{450}{nm}}$ ; solution 2 : $\res{\lambda_{\max}\approx\qty{525}{nm}}$.
\item La solution 1 absorbe surtout le \res{\text{bleu}}, la solution 2 surtout le \res{\text{vert}}.
\item La solution 1 paraît \res{\text{jaune-orangé}}, la solution 2 \res{\text{magenta}}.
\item Le permanganate de potassium en solution est \res{\text{violet-rose}} (magenta).
\end{questions}
\schema[corrige]{02-ex01}
\end{corrige}
\end{blocLivre}
\begin{blocLivre}{cor-2-2}
\begin{corrige}{2}{Prévoir l'allure d'un spectre}
\demarche{La couleur absorbée est la complémentaire de la couleur perçue.}
\begin{questions}
\item Elle absorbe la complémentaire du bleu-cyan : le \res{\text{rouge-orangé}}, au-delà de \qty{600}{nm} environ.
\item Elle absorbe le \res{\text{bleu-violet}}, vers 400 à \qty{480}{nm}.
\item Elle absorbe le bleu et le rouge et transmet le vert : un maximum vers \res{\qty{430}{nm}} et un autre vers \res{\qty{660}{nm}}.
\item Elle est \res{\text{incolore}} ; on ne peut pas la doser dans le visible (mais parfois dans l'ultraviolet).
\end{questions}
\end{corrige}
\end{blocLivre}
\begin{blocLivre}{cor-2-3}
\begin{corrige}{3}{Appliquer la loi de Beer-Lambert}
\demarche{$A=\varepsilon\,\ell\,c$, donc $c=\dfrac{A}{\varepsilon\,\ell}$ ; puis $t=c\times M$.}
\begin{questions}
\item $A=\num{2.4e3}\times1{,}00\times\num{2.0e-4}=\res{0{,}48}$.
\item $c=\dfrac{0{,}30}{\num{2.4e3}\times1{,}00}=\res{\qty{1,25e-4}{mol.L^{-1}}}$.
\item $A$ est proportionnelle à $\ell$ : elle double et vaut \res{0{,}96}.
\item $M(\ce{KMnO4})=39{,}1+54{,}9+4\times16{,}0=\qty{158,0}{g.mol^{-1}}$ et $t=\num{1.25e-4}\times158{,}0=\qty{1,98e-2}{g.L^{-1}}$, soit \res{\qty{19,8}{mg.L^{-1}}}.
\end{questions}
\piege{Mélanger les unités : avec $\varepsilon$ en \unit{L.mol^{-1}.cm^{-1}}, la longueur $\ell$ doit être en centimètres et $c$ s'obtient en \unit{mol.L^{-1}}.
L'absorbance, elle, n'a pas d'unité.}
\end{corrige}
\end{blocLivre}
\begin{blocLivre}{cor-2-4}
\begin{corrige}{4}{Lire une droite d'étalonnage}
\demarche{Une droite passant par l'origine traduit $A=k\,c$ ; $k$ est son coefficient directeur, et $\varepsilon=\dfrac{k}{\ell}$.}
\begin{questions}
\item Les points sont alignés sur une \res{\text{droite passant par l'origine}} : $A$ est proportionnelle à $c$, la loi est vérifiée.
\item $k=\dfrac{0{,}96}{80}=\res{\qty{0,012}{L.mmol^{-1}}}$, soit \qty{12}{L.mol^{-1}} pour une cuve de \qty{1,00}{cm} : $\res{\varepsilon=\qty{12}{L.mol^{-1}.cm^{-1}}}$.
\item $c=\dfrac{0{,}54}{0{,}012}=\res{\qty{45}{mmol.L^{-1}}}$.
\item $t=0{,}045\times159{,}6=\res{\qty{7,18}{g.L^{-1}}}$.
\end{questions}
\schema[corrige]{02-ex04}
\end{corrige}
\end{blocLivre}
\begin{blocLivre}{cor-2-5}
\begin{corrige}{5}{Diluer avant de mesurer}
\demarche{La loi n'est fiable que pour $A<2$ environ ; on mesure sur la solution diluée, puis on multiplie par le facteur de dilution.}
\begin{questions}
\item Une absorbance de 3,1 sort du domaine où la loi de Beer-Lambert s'applique : le résultat \res{\text{n'est pas fiable}}.
\item $c=\dfrac{0{,}42}{0{,}106}=\res{\qty{3,96}{\micro mol.L^{-1}}}$ dans le sirop dilué, donc $\res{\qty{39,6}{\micro mol.L^{-1}}}$ dans le sirop.
\item Prélever \qty{10,0}{mL} de sirop avec une pipette jaugée, les verser dans une \res{\text{fiole jaugée de 100,0 mL}}, compléter avec de l'eau
distillée jusqu'au trait de jauge, boucher et homogénéiser.
\end{questions}
\schema[corrige]{02-ex05-cor}
\retiens{On mesure sur la solution diluée, puis on revient à la solution de départ en multipliant par le facteur de dilution.}
\end{corrige}
\end{blocLivre}
\begin{blocLivre}{cor-2-6}
\begin{corrige}{6}{Préparer une gamme étalon}
\demarche{La quantité de soluté se conserve lors d'une dilution : $c_0\,V_0=c\,V$.}
\begin{questions}
\item $V_0=\dfrac{c\,V}{c_0}$, avec $V=\qty{25,0}{mL}$.
\item Pour 1,0 ; 2,0 ; 4,0 ; 6,0 et \qty{8,0}{\micro mol.L^{-1}} : $V_0=\res{2{,}5}$ ; \res{5{,}0} ; \res{10{,}0} ; \res{15{,}0} et \res{\qty{20,0}{mL}}.
\item Une \res{\text{burette graduée}} permet de verser tous ces volumes avec précision (des pipettes jaugées ne conviendraient
que pour 5,0 ; 10,0 et \qty{20,0}{mL}).
\item Le coefficient $k=\varepsilon\,\ell$ dépend de la longueur d'onde et de l'épaisseur de la cuve : il doit être le \res{\text{même}} pour toutes les
mesures, sinon les points ne sont pas comparables.
\end{questions}
\end{corrige}
\end{blocLivre}
\begin{blocLivre}{cor-2-7}
\begin{corrige}{7}{Le permanganate de l'eau de Dakin}
\demarche{$c=\dfrac{A}{\varepsilon\,\ell}$, puis $t=c\times M$ ; comparer à l'étiquette convertie dans la même unité.}
\begin{questions}
\item Le permanganate absorbe le vert : la solution paraît \res{\text{rose}}, couleur complémentaire. On mesure au maximum d'absorption,
\qty{525}{nm}, où la mesure est la plus sensible.
\item Non : l'hypochlorite est \res{\text{incolore}}, il n'absorbe pas à \qty{525}{nm}.
\item $c=\dfrac{0{,}150}{\num{2.4e3}\times1{,}00}=\res{\qty{6,25e-5}{mol.L^{-1}}}$, d'où $t=\num{6.25e-5}\times158{,}0=\qty{9,9e-3}{g.L^{-1}}=\res{\qty{9,9}{mg.L^{-1}}}$.
\item \qty{0,0010}{g} pour \qty{100}{mL} correspondent à \qty{0,010}{g.L^{-1}}, soit \res{\qty{10}{mg.L^{-1}}} : la mesure est \res{\text{conforme}} à l'étiquette.
\end{questions}
\retiens{On règle le spectrophotomètre à la longueur d'onde du maximum d'absorption : l'absorbance y est la plus grande,
donc la mesure la plus sensible et la moins affectée par un petit décalage de longueur d'onde.}
\end{corrige}
\end{blocLivre}
\begin{blocLivre}{cor-2-8}
\begin{corrige}{8}{Le cuivre d'une bouillie bordelaise}
\demarche{$c=\dfrac{A}{k}$, puis $n=c\times V$ et $m=n\times M$ ; pourcentage $=\dfrac{m(\ce{Cu})}{m(\text{produit})}$.}
\begin{questions}
\item $c=\dfrac{0{,}37}{0{,}012}=\res{\qty{30,8}{mmol.L^{-1}}}$.
\item $n=\num{30.8e-3}\times0{,}100=\qty{3,08e-3}{mol}$ et $m=\num{3.08e-3}\times63{,}5=\res{\qty{0,196}{g}}$.
\item $\dfrac{0{,}196}{1{,}00}=\res{\qty{19,6}{\percent}}$ : à la précision de la mesure, le produit est \res{\text{conforme}} aux 20~\% annoncés.
\item La droite n'a été vérifiée que sur le domaine de la gamme (jusqu'à $A=0{,}96$) : en dehors, rien ne garantit la
proportionnalité.
\end{questions}
\end{corrige}
\end{blocLivre}
\begin{blocLivre}{cor-2-9}
\begin{corrige}{9}{Tester la loi de Beer-Lambert en Python}
\demarche{$k=\dfrac{\sum c_i\,A_i}{\sum c_i^2}$ est le coefficient de la droite passant par l'origine la plus proche des points.}
\begin{questions}
\item 0.1061 est le coefficient $\res{k}$ de la relation $A=k\,c$, en \res{\unit{L.\micro mol^{-1}}} puisque $c$ est en \unit{\micro mol.L^{-1}}.
\item Elle calcule, pour chaque point, l'\res{\text{écart relatif}} entre l'absorbance mesurée et celle de la droite, en pourcentage. Les
écarts restent inférieurs à 1~\% en valeur absolue.
\item Oui : les points sont tous très proches d'une \res{\text{droite passant par l'origine}}.
\item \code{print(0.50 / k)}, qui affiche environ 4,71 : $\res{c\approx\qty{4,71}{\micro mol.L^{-1}}}$.
\end{questions}
\end{corrige}
\end{blocLivre}
\begin{blocLivre}{cor-2-10}
\begin{corrige}{10}{Les limites de la loi de Beer-Lambert}
\demarche{La loi n'est valable que là où les points restent sur la droite passant par l'origine.}
\begin{questions}
\item Les points de 5 à \qty{15}{\micro mol.L^{-1}} sont sur la droite $A=0{,}106\,c$ ; au-delà, $A$ augmente moins vite. La loi est vérifiée
jusqu'à \res{\qty{15}{\micro mol.L^{-1}}} environ, soit $A$ inférieure à 1,6.
\item $0{,}106\times25=\res{2{,}65}$, alors que la mesure donne 2,31 : la loi ne s'applique plus.
\item \res{\text{Non}} : cette absorbance est hors du domaine de validité. On dilue la solution quatre fois, par exemple \qty{25,0}{mL}
dans une fiole jaugée de \qty{100,0}{mL}, pour obtenir une absorbance d'environ 0,6.
\end{questions}
\schema[corrige]{02-ex10}
\retiens{La loi de Beer-Lambert ne vaut que pour des solutions assez diluées ; en pratique, on travaille avec des absorbances
inférieures à 1,5 environ, en vérifiant que les étalons sont alignés avec l'origine.}
\end{corrige}
\end{blocLivre}
\begin{blocLivre}{cor-2-11}
\begin{corrige}{11}{Un sirop vert, deux colorants}
\demarche{Le sirop absorbe le bleu et le rouge-orangé : il transmet surtout le vert.}
\begin{questions}
\item Le sirop absorbe vers \qty{427}{nm} (bleu) et vers \qty{638}{nm} (rouge-orangé) : la lumière transmise est surtout \res{\text{verte}}.
\item La tartrazine, jaune, absorbe le bleu : maximum à \res{\qty{427}{nm}} ; le bleu patenté absorbe le rouge-orangé : maximum à \res{\qty{638}{nm}}.
\item À \qty{638}{nm}, seul le bleu patenté absorbe : l'absorbance du sirop est \res{\text{celle du bleu patenté seul}}.
\item $c=\dfrac{0{,}53}{0{,}106}=\res{\qty{5,0}{\micro mol.L^{-1}}}$.
\end{questions}
\schema[corrige]{02-ex11}
\end{corrige}
\end{blocLivre}
\begin{blocLivre}{cor-2-12}
\begin{corrige}{12}{Combien de verres de sirop ?}
\demarche{$c=\dfrac{A}{k}$, facteur de dilution 20, $t=c\times M$ ; masse par verre, puis comparaison à la DJA.}
\begin{questions}
\item $c=\dfrac{0{,}53}{0{,}106}=\qty{5,0}{\micro mol.L^{-1}}$ dans la solution diluée, donc $c_\text{sirop}=20\times5{,}0=\res{\qty{100}{\micro mol.L^{-1}}}$.
\item $t=\num{100e-6}\times582{,}7=\qty{0,058}{g.L^{-1}}$, soit \res{\qty{58}{mg.L^{-1}}}.
\item $m=58\times\num{0.020}\approx\res{\qty{1,2}{mg}}$ par verre.
\item La DJA d'un enfant de \qty{30}{kg} vaut $5\times30=\qty{150}{mg}$ par jour, soit près de \res{\text{130 verres}} : aucun risque de la dépasser
avec une consommation normale.
\end{questions}
\end{corrige}
\end{blocLivre}
\begin{blocLivre}{cor-2-13}
\begin{corrige}{13}{Les nitrates d'une eau de puits}
\demarche{Coefficient de la droite d'étalonnage, puis $t=\dfrac{A}{k}$ et $c=\dfrac{t}{M}$ ; comparer à la limite de \qty{50}{mg.L^{-1}}.}
\begin{questions}
\item Les ions nitrate sont incolores : ils n'absorbent pas dans le visible. Le réactif les transforme en une \res{\text{espèce colorée}}
dont l'absorbance est proportionnelle à leur concentration.
\item $\dfrac{A}{t}$ vaut toujours $0{,}0085$ : la loi est vérifiée, avec $\res{k=\qty{0,0085}{L.mg^{-1}}}$.
\item $t=\dfrac{0{,}392}{0{,}0085}=\res{\qty{46,1}{mg.L^{-1}}}$ et, avec $M(\ce{NO3-})=\qty{62,0}{g.mol^{-1}}$, $c=\res{\qty{7,44e-4}{mol.L^{-1}}}$.
\item Puits A : \qty{46,1}{mg.L^{-1}}, inférieur à \qty{50}{mg.L^{-1}} : \res{\text{potable}}, mais de peu. Puits B : $A=0{,}48$ dépasse la gamme
($A\leq0{,}425$) ; en extrapolant, $t\approx\qty{56}{mg.L^{-1}}$, au-dessus de la limite : \res{\text{non potable}}. Pour confirmer, il faudrait
diluer l'eau du puits B ou prolonger la gamme.
\end{questions}
\schema[corrige]{02-ex13-cor}
\piege{Lire une absorbance hors de la gamme d'étalonnage en prolongeant la droite : rien ne garantit que la loi reste
valable. On dilue l'échantillon pour revenir dans la gamme, puis on multiplie par le facteur de dilution.}
\end{corrige}
\end{blocLivre}
\begin{blocLivre}{cor-2-14}
\begin{corrige}{14}{La caféine d'une boisson énergisante}
\demarche{$A=\varepsilon\,\ell\,c$, facteur de dilution 50, $t=c\times M$.}
\begin{questions}
\item $c=\dfrac{0{,}32}{194{,}2}=\qty{1,65e-3}{mol.L^{-1}}$, d'où $A=\num{9.7e3}\times1{,}00\times\num{1.65e-3}\approx\res{16}$ : bien au-delà du domaine de
validité de la loi, d'où la dilution.
\item $c=\dfrac{0{,}318}{\num{9.7e3}}=\res{\qty{3,28e-5}{mol.L^{-1}}}$ dans la solution diluée, donc $\res{\qty{1,64e-3}{mol.L^{-1}}}$ dans la boisson.
\item $t=\num{1.64e-3}\times194{,}2=\res{\qty{0,318}{g.L^{-1}}}$, soit \qty{31,8}{mg} pour \qty{100}{mL} : \res{\text{conforme}} à l'étiquette.
\item Limite : $3\times60=\qty{180}{mg}$ par jour ; une canette apporte $0{,}318\times250=\qty{80}{mg}$ : au plus \res{\text{deux canettes}} par jour
($\dfrac{180}{80}=2{,}25$), sans compter les autres sources de caféine.
\end{questions}
\end{corrige}
\end{blocLivre}
\begin{blocLivre}{cor-2-15}
\begin{corrige}{15}{Le fer d'un complément alimentaire}
\demarche{$c=\dfrac{A}{\varepsilon\,\ell}$, facteur de dilution 25, puis $m=c\times V\times M$ pour la fiole de \qty{1,000}{L}.}
\begin{questions}
\item $F=\dfrac{50{,}0}{2{,}00}=\res{25}$.
\item $c=\dfrac{0{,}108}{\num{1.11e4}\times1{,}00}=\res{\qty{9,73e-6}{mol.L^{-1}}}$.
\item Dans la fiole de \qty{1,000}{L} : $25\times\num{9.73e-6}=\qty{2,43e-4}{mol.L^{-1}}$, soit $n=\qty{2,43e-4}{mol}$ et $m=\num{2.43e-4}\times55{,}8=\res{\qty{13,6}{mg}}$ de fer :
97~\% de l'indication, ce qui est \res{\text{conforme}} compte tenu de la précision de la mesure.
\item $c=\dfrac{\num{14e-3}}{55{,}8\times25}=\qty{1,00e-5}{mol.L^{-1}}$ dans la fiole de \qty{50,0}{mL}, d'où $A=\res{0{,}111}$.
\end{questions}
\end{corrige}
\end{blocLivre}
